\documentclass[11pt,a4paper]{article}

% ---------- Encoding e idioma ----------
\usepackage[utf8]{inputenc}
\usepackage[T1]{fontenc}
\usepackage[spanish,es-tabla]{babel}

% ---------- Diseño de página ----------
\usepackage[margin=2.5cm]{geometry}
\usepackage{setspace}
\onehalfspacing

% ---------- Matemáticas ----------
\usepackage{amsmath,amssymb,amsthm,mathtools}

% ---------- Figuras y tablas ----------
\usepackage{graphicx}
\usepackage{booktabs}
\usepackage{caption}
\usepackage{float}

% ---------- Color y cajas ----------
\usepackage{xcolor}
\usepackage{tcolorbox}
\tcbuselibrary{skins, breakable}

\definecolor{colLAB}{RGB}{155, 89, 182}
\definecolor{colCLAHE}{RGB}{52, 152, 219}
\definecolor{colTILES}{RGB}{230, 126, 34}
\definecolor{colCLIP}{RGB}{231, 76, 60}
\definecolor{colCLAMP}{RGB}{39, 174, 96}
\definecolor{colGRAY}{RGB}{127, 140, 141}
\definecolor{colMED}{RGB}{142, 68, 173}

\newtcolorbox{analogy}[1][]{
    colback=yellow!10,
    colframe=yellow!60!orange,
    boxrule=1pt,
    arc=3pt,
    title=\textbf{Analogía},
    #1
}

\newtcolorbox{codebox}[1][]{
    colback=gray!10,
    colframe=gray!50,
    boxrule=0.5pt,
    arc=2pt,
    fontupper=\ttfamily\small,
    #1
}

\newtcolorbox{keybox}[2][]{
    colback=#2!8,
    colframe=#2!80,
    boxrule=1.2pt,
    arc=4pt,
    title={\Large\textbf{#1}},
    coltitle=white,
    colbacktitle=#2,
    breakable,
    fonttitle=\sffamily
}

% ---------- Código ----------
\usepackage{listings}
\lstdefinestyle{matlab}{
    language=Matlab,
    backgroundcolor=\color{gray!8},
    basicstyle=\ttfamily\small,
    keywordstyle=\color{blue!70!black}\bfseries,
    commentstyle=\color{green!50!black}\itshape,
    stringstyle=\color{red!60!black},
    numbers=left,
    numberstyle=\tiny\color{gray},
    frame=single,
    rulecolor=\color{black!30},
    breaklines=true,
    showstringspaces=false,
    tabsize=2,
    captionpos=b
}
\lstset{style=matlab}

% ---------- Enlaces ----------
\usepackage[colorlinks=true, linkcolor=blue!60!black, urlcolor=blue!60!black]{hyperref}

% ---------- Metadatos ----------
\title{
    \vspace{-1.5cm}
    \textbf{Conceptos Clave de Procesamiento de Imágenes} \\[0.3cm]
    \Large Glosario visual: CLAHE, LAB, Tiles, Clip Limit y Clamping
}
\author{Jamin Yauri Cajas}
\date{\today}

% ============================================================
\begin{document}
\maketitle

\begin{abstract}
\noindent
Documento de apoyo al informe principal del proyecto de Visión por Computadora.
Explica, con analogías simples, ejemplos visuales y fórmulas, los cinco conceptos
técnicos más importantes detrás del pipeline de preprocesamiento: el espacio de
color \textbf{Lab}, la ecualización adaptativa con restricción de contraste
\textbf{CLAHE}, la idea de \textbf{tiles} o baldosas, el mecanismo de
\textbf{clip limit} y la operación de \textbf{clamping}. Se incluyen además los
conceptos complementarios del \textbf{filtro de mediana} y el \textbf{balance de
blancos gray-world}.
\end{abstract}

\tableofcontents
\vspace{1cm}

% ============================================================
\section{Motivación}

Un pipeline de preprocesamiento, por más corto que sea, combina varios operadores
de procesamiento de imágenes con funciones distintas. Entender los nombres técnicos
que aparecen en el código (\texttt{rgb2lab}, \texttt{adapthisteq},
\texttt{ClipLimit}, \texttt{NumTiles}, \texttt{medfilt2}, \texttt{max/min} como
clamping) permite justificar cada decisión del pipeline y responder con solvencia
a las preguntas del enunciado. Este documento presenta esos conceptos uno por
uno, con la misma profundidad que se usaría en una clase del curso.

% ============================================================
\section{Espacio de color \textcolor{colLAB}{Lab}}
\begin{keybox}[LAB — Separar brillo del color]{colLAB}

\subsection*{¿Qué es?}

Una forma distinta de representar colores donde el \textbf{brillo} y el
\textbf{color puro} se guardan en canales separados.

\subsection*{Comparación con RGB}

\textbf{RGB} almacena tres intensidades —Rojo, Verde, Azul— que se suman para
formar el color final. Cualquier operación que cambia solo una intensidad
(subir el rojo) también altera el brillo y la saturación simultáneamente. No
hay forma de decir \emph{``solo quiero oscurecer este píxel sin tocar su tono''}.

\textbf{Lab} separa explícitamente:
\begin{itemize}
    \item \(L\) (\emph{Lightness}) — canal de brillo: \(0\) es negro puro, \(100\) es blanco puro.
    \item \(a\) — eje cromático verde \(\leftrightarrow\) rojo.
    \item \(b\) — eje cromático azul \(\leftrightarrow\) amarillo.
\end{itemize}

La transformación de RGB a Lab no es lineal: pasa por un espacio intermedio
(XYZ de la CIE) y aplica funciones de cube-root diseñadas para imitar la
respuesta no lineal del ojo humano a la luz.

\begin{analogy}
Piensa en los controles \emph{Brillo} y \emph{Tono} de un televisor. En RGB es
imposible tocar uno sin tocar el otro; en Lab son perillas independientes.
Cambiar \(L\) es como mover solo la perilla de brillo.
\end{analogy}

\subsection*{¿Por qué se usa en nuestro pipeline?}

Para aplicar CLAHE solo sobre el brillo (canal \(L\)) sin modificar el tono
(canales \(a, b\)). Si aplicáramos CLAHE en RGB píxel a píxel, cada canal se
redistribuiría de forma desigual, desplazando el color percibido y generando
\emph{aberraciones cromáticas} (un rojo volviéndose anaranjado, por ejemplo).

\subsection*{Uso en el código}

\begin{codebox}
lab = rgb2lab(img);       \% RGB $\to$ Lab \\
L = lab(:,:,1) / 100;     \% solo el canal de brillo \\
\% ... aquí se aplica CLAHE a L ... \\
lab(:,:,1) = L * 100;     \% se vuelve a insertar \\
img = lab2rgb(lab);       \% Lab $\to$ RGB
\end{codebox}

\subsection*{Alternativas}

Otros espacios que también separan brillo de color son \textbf{HSV} (Hue,
Saturation, Value) y \textbf{YCbCr}. Se prefiere Lab porque es
\emph{perceptualmente uniforme}: distancias iguales en el espacio Lab
corresponden a diferencias iguales percibidas por el ojo, lo que no ocurre en
HSV o YCbCr.

\end{keybox}

% ============================================================
\section{\textcolor{colCLAHE}{CLAHE} — Ecualización Adaptativa con Límite de Contraste}
\begin{keybox}[CLAHE — Mejorar contraste zona por zona]{colCLAHE}

\subsection*{¿Qué es?}

Un algoritmo para aumentar el contraste local de una imagen, adaptándose a
las condiciones de cada región \textbf{independientemente} del resto.

\subsection*{El problema que resuelve}

La \textbf{ecualización de histograma clásica} busca que todos los tonos de
gris aparezcan con la misma frecuencia. Tiene dos problemas graves:

\begin{enumerate}
    \item \textbf{No es adaptativa.} Aplica una única transformación a toda
          la imagen. Si una zona está en sombra y otra bien iluminada, la
          transformación no puede beneficiar a ambas a la vez.
    \item \textbf{Amplifica ruido.} En zonas planas (cielo, pared blanca), el
          histograma es muy estrecho; al ``estirarlo'' se magnifican las
          micro-variaciones del sensor.
\end{enumerate}

CLAHE resuelve ambos:
\begin{itemize}
    \item \textbf{Adaptive}: divide la imagen en bloques y ecualiza cada uno
          por separado, ajustándose al histograma local.
    \item \textbf{Contrast Limited}: recorta los picos del histograma antes de
          ecualizar, evitando amplificar ruido.
\end{itemize}

\begin{analogy}
Una foto con una zona oscura en sombra y otra quemada por el sol es como un
auditorio donde una mitad está a oscuras y la otra con los focos encendidos.
Subir el brillo general no ayuda a nadie. CLAHE es como ponerle un dimmer
independiente a cada fila de butacas: cada una recibe exactamente la luz
que necesita.
\end{analogy}

\subsection*{Algoritmo paso a paso}

\begin{enumerate}
    \item Dividir la imagen en una grilla de tiles (típicamente \(8\times 8\)).
    \item Para cada tile, calcular su histograma \(h_T(k)\).
    \item Aplicar el \emph{clip limit} \(\beta\) al histograma: los bins que
          superan el umbral se recortan, y el excedente se redistribuye uniformemente.
    \item Calcular la CDF (función de distribución acumulada) del histograma
          recortado; esa es la función de mapeo que ecualiza el tile.
    \item Aplicar la transformación \textbf{interpolando bilinealmente} entre
          tiles vecinos para que no se note la cuadrícula.
\end{enumerate}

\subsection*{Uso en el código}

\begin{codebox}
L = adapthisteq(L, ... \\
\quad 'ClipLimit',    0.005, ... \\
\quad 'NumTiles',     [8 8], ... \\
\quad 'Distribution', 'uniform');
\end{codebox}

\end{keybox}

% ============================================================
\section{\textcolor{colTILES}{Tiles} — Baldosas para procesamiento local}
\begin{keybox}[TILES — Dividir para conquistar]{colTILES}

\subsection*{¿Qué son?}

Los \textbf{tiles} (``baldosas'' o ``mosaicos'' en español) son las regiones
rectangulares en las que CLAHE divide la imagen para procesarla por partes.
La grilla típica es de \(8 \times 8 = 64\) tiles.

\subsection*{Representación visual}

\begin{center}
\begin{tabular}{|c|c|c|c|c|c|c|c|}
\hline
T$_{1,1}$ & T$_{1,2}$ & T$_{1,3}$ & T$_{1,4}$ & T$_{1,5}$ & T$_{1,6}$ & T$_{1,7}$ & T$_{1,8}$ \\
\hline
T$_{2,1}$ & T$_{2,2}$ & T$_{2,3}$ & T$_{2,4}$ & T$_{2,5}$ & T$_{2,6}$ & T$_{2,7}$ & T$_{2,8}$ \\
\hline
\multicolumn{8}{|c|}{$\cdots$ hasta completar 8 filas $\cdots$} \\
\hline
T$_{8,1}$ & T$_{8,2}$ & T$_{8,3}$ & T$_{8,4}$ & T$_{8,5}$ & T$_{8,6}$ & T$_{8,7}$ & T$_{8,8}$ \\
\hline
\end{tabular}
\end{center}

Cada tile tiene su propio histograma y su propia función de ecualización.

\subsection*{¿Por qué 8×8?}

Es un compromiso:

\begin{table}[H]
\centering
\begin{tabular}{@{}lll@{}}
\toprule
Tamaño de grilla & Tile grande/pequeño & Problema \\
\midrule
\(2 \times 2\)   & Muy grandes & Se comporta como ecualización global \\
\(8 \times 8\)   & \textbf{Equilibrado} & \textbf{Estándar de la literatura} \\
\(32 \times 32\) & Muy pequeños & Histograma no representativo (ruido) \\
\bottomrule
\end{tabular}
\end{table}

\subsection*{El truco de la interpolación bilineal}

Si se aplicara cada transformación solo a su tile, se verían bordes visibles
entre tiles vecinos (\emph{blocking artifacts}). CLAHE interpola
bilinealmente: cada píxel recibe un promedio ponderado de las transformaciones
de los cuatro tiles más cercanos, en función de su distancia a los centros de
cada uno. El resultado es una transición continua.

\end{keybox}

% ============================================================
\section{\textcolor{colCLIP}{Clip Limit} — Controlar la amplificación de ruido}
\begin{keybox}[CLIP LIMIT — El freno anti-ruido]{colCLIP}

\subsection*{¿Qué es?}

El \textbf{clip limit} es el tope máximo que puede tener cualquier bin del
histograma dentro de un tile antes de ser recortado. Controla cuán agresiva
puede ser la ecualización.

\subsection*{Por qué existe}

Sin clip limit, CLAHE amplificaría ferozmente las zonas planas. Supongamos
que un tile representa la madera del escritorio, casi uniforme: su histograma
tendría un pico altísimo en un solo tono y nada en el resto. Ecualizar eso
significa ``estirar'' ese pico sobre todo el rango, lo que magnifica las
micro-variaciones del sensor (ruido) y las hace visibles.

\begin{analogy}
Imagina un escáner médico subido al máximo. Revela tejidos finísimos, pero
también amplifica el ruido eléctrico: aparecen granos falsos. El clip limit
es el nivel de ganancia máxima antes de que empiecen a aparecer granos.
\end{analogy}

\subsection*{Representación visual del clipping}

\begin{verbatim}
    Histograma original       Histograma clipeado al umbral beta·N

         |                          |
         |  #                       |
         |  #                       |  ===== <- umbral beta*N
         |  #                       |  # # #
         |  #                       |  # # #
         |  # #                     |  # # #
         |  # # #                   |  # # # #
         |  # # # #                 |  # # # # #
         +----------------+         +----------------+
         0  valor 1  ... 255        0  valor 1  ... 255
\end{verbatim}

El ``excedente'' (la parte por encima del umbral) no se descarta: se
redistribuye uniformemente entre todos los bins para conservar la masa total
del histograma.

\subsection*{Fórmula}

Para un tile \(T\) con \(N\) píxeles totales y límite relativo \(\beta\):
\begin{equation}
    h_{\text{clip}}(k) = \min\!\bigl(h_T(k),\; \beta N\bigr)
\end{equation}
El excedente \(\sum_k \bigl[h_T(k) - h_{\text{clip}}(k)\bigr]\) se suma
equitativamente a todos los bins.

\subsection*{Valor usado en el código}

En nuestro pipeline se usa \(\beta = 0{,}005\) (0.5\% de los píxeles por bin).
Es un valor \textbf{conservador}: da realce sin introducir ruido visible.
Valores mayores (\(0{,}02\) o más) producen imágenes muy contrastadas pero
ruidosas.

\end{keybox}

% ============================================================
\section{\textcolor{colCLAMP}{Clamping} — Cortar al rango válido}
\begin{keybox}[CLAMPING — Mantener los píxeles en rango]{colCLAMP}

\subsection*{¿Qué es?}

Una operación que fuerza a cada valor a quedarse dentro de un rango permitido.
Cualquier valor por debajo del mínimo se pega al mínimo, y cualquier valor
por encima del máximo se pega al máximo.

\subsection*{Fórmula}

\begin{equation}
    \text{clamp}(x; a, b) =
    \begin{cases}
        a & \text{si } x < a, \\
        x & \text{si } a \le x \le b, \\
        b & \text{si } x > b.
    \end{cases}
\end{equation}

En forma compacta usando \(\min\) y \(\max\):
\begin{equation}
    \text{clamp}(x; a, b) = \max\!\bigl(a,\; \min(b, x)\bigr).
\end{equation}

En nuestro pipeline \(a = 0\) y \(b = 1\) (imágenes en \texttt{double}).

\begin{analogy}
El control de volumen del celular: va de 0 a 100. Presionar ``subir'' cuando
ya estás en 100 no pasa a 101 — se queda en 100. Presionar ``bajar'' cuando
estás en 0 no pasa a \(-1\) — se queda en 0.
\end{analogy}

\subsection*{¿Por qué hace falta?}

La conversión \(\text{Lab} \to \text{RGB}\) tras CLAHE puede producir valores
ligeramente fuera del rango \([0, 1]\) debido a:
\begin{itemize}
    \item Errores de redondeo en aritmética de punto flotante.
    \item Puntos del espacio Lab que están fuera del gamut visible de RGB (colores
          que existen en Lab pero no pueden representarse en RGB).
\end{itemize}

Si no se clampa, al guardar la imagen como \texttt{uint8} MATLAB aplica
truncamiento o wrap-around, lo que produce ``píxeles chillones'' (rojo puro
donde debería haber blanco, por ejemplo). Clamping evita ese artefacto.

\subsection*{Uso en el código}

\begin{codebox}
img = max(0, min(1, img));
\end{codebox}

\end{keybox}

% ============================================================
\section{\textcolor{colMED}{Filtro de mediana} — Denoise que respeta bordes}
\begin{keybox}[MEDIANA — Filtro no lineal anti-impulso]{colMED}

\subsection*{¿Qué hace?}

Reemplaza cada píxel por la \textbf{mediana} (valor central ordenado) de los
píxeles de su vecindad.

\subsection*{Mediana vs. promedio}

Sea la vecindad \(\{10, 12, 11, 255, 10, 13, 11, 12, 10\}\) (ruido sal):
\begin{itemize}
    \item \textbf{Promedio} = \(\frac{10+12+11+255+10+13+11+12+10}{9} \approx 38\) — contaminado por el outlier.
    \item \textbf{Mediana} = \(11\) — el outlier se descarta.
\end{itemize}

\subsection*{Propiedades}

\begin{enumerate}
    \item \textbf{No lineal}: no se puede expresar como convolución con un kernel.
    \item \textbf{Robusto a outliers}: píxeles aislados extremos no afectan la salida.
    \item \textbf{Preserva bordes} mejor que un filtro promedio, porque no introduce valores intermedios inventados.
    \item \textbf{Ideal contra ruido impulsivo} (sal y pimienta); poco efectivo contra ruido gaussiano — por eso se complementa con filtro gaussiano.
\end{enumerate}

\subsection*{Uso en el código}

\begin{codebox}
img(:,:,c) = medfilt2(img(:,:,c), [3 3]);
\end{codebox}

\end{keybox}

% ============================================================
\section{\textcolor{colGRAY}{Gray-world} — Balance de blancos automático}
\begin{keybox}[GRAY-WORLD — Corregir el tinte de iluminación]{colGRAY}

\subsection*{La hipótesis}

En una escena natural típica, la suma de todos los colores presentes se
aproxima a gris neutro. Si en una foto el promedio del canal rojo es
significativamente mayor que el del verde o azul, es porque la
\textbf{iluminación} tiene un tinte cálido (bombilla incandescente), no
porque la escena sea ``roja''.

\subsection*{La corrección}

Se calcula la media de cada canal y una media global, y se escala cada canal
para que su media iguale la global:
\begin{equation}
    \bar{c} = \frac{1}{HW}\sum_{x,y} I(x,y,c), \qquad
    \bar{m} = \frac{\bar{R} + \bar{G} + \bar{B}}{3},
\end{equation}
\begin{equation}
    I'(x,y,c) = I(x,y,c) \cdot \frac{\bar{m}}{\bar{c}}.
\end{equation}

\subsection*{Ventajas y limitaciones}

\textbf{Ventajas:}
\begin{itemize}
    \item No tiene parámetros; se auto-calibra a partir de cada imagen.
    \item Muy rápido.
    \item Funciona bien en escenas con variedad cromática.
\end{itemize}

\textbf{Limitaciones:}
\begin{itemize}
    \item Falla si la escena está dominada por un color (p.\,ej., solo
          marcadores rojos): creería que la iluminación es cálida y
          oscurecería el rojo.
    \item Alternativas más robustas: \emph{Shades of Gray}, \emph{max-RGB},
          \emph{Gray-Edge}, todas también automáticas.
\end{itemize}

\end{keybox}

% ============================================================
\section{Resumen integrado}

La Tabla~\ref{tab:resumen} condensa la función, el tipo y el valor usado de
cada concepto revisado.

\begin{table}[H]
\centering
\caption{Resumen de los conceptos clave.}
\label{tab:resumen}
\renewcommand{\arraystretch}{1.3}
\begin{tabular}{@{}llll@{}}
\toprule
Concepto & Tipo & Función & Valor/Parámetro \\
\midrule
Lab                 & Espacio de color      & Separa brillo de color       & --- \\
Mediana             & Filtro no lineal      & Elimina outliers             & kernel \([3,3]\) \\
Gaussiano           & Filtro lineal         & Suaviza ruido residual       & \(\sigma = 0{,}8\) \\
Gray-world          & Algoritmo de WB       & Corrige tinte de iluminación & sin parámetros \\
CLAHE               & Realce de contraste   & Ecualiza localmente          & ver abajo \\
Tiles               & Parámetro de CLAHE    & Grilla de división           & \(8\times 8\) \\
Clip limit          & Parámetro de CLAHE    & Controla amplificación ruido & \(\beta = 0{,}005\) \\
Clamping            & Operación aritmética  & Fuerza rango \([0,1]\)       & \(\max(0, \min(1, x))\) \\
\bottomrule
\end{tabular}
\end{table}

% ============================================================
\section{Conclusión}

Los cinco conceptos principales — \textbf{Lab}, \textbf{CLAHE}, \textbf{tiles},
\textbf{clip limit} y \textbf{clamping} — no son términos aislados sino piezas
de un mismo razonamiento: se quiere mejorar el contraste de la imagen \emph{sin
alterar su color} (de ahí Lab), \emph{adaptándose a cada región} (de ahí los
tiles que alimentan CLAHE), \emph{sin amplificar el ruido de fondo} (de ahí el
clip limit), \emph{y garantizando que el resultado sea una imagen válida} (de
ahí el clamping final). Los dos conceptos complementarios —
\textbf{filtro de mediana} y \textbf{gray-world} — completan el pipeline
atacando los dos defectos restantes: ruido impulsivo y dominante cromática
de iluminación. El conjunto forma un algoritmo automático, adaptativo y
justificable etapa por etapa.

\end{document}
